\section{Mixture-of-experts workloads}
\label{sec:moe}

\begin{question}
  \textbf{How does GLM-5 behave at scale?} What does expert traffic cost when it leaves the node or crosses
  switches, and what does FP8 bring?
\end{question}

Expert parallelism (EP) adds all-to-all exchanges every MoE layer. With EP$\le$4, experts stay inside a
node (NVLink); beyond that, all-to-all traffic crosses the single HCA of each node and, at larger EP, leaf
switches. All-to-all bandwidth across leaves is the weakest scope measured (\cref{fig:comm}).

\begin{figure}[h]
  \centering\fig{moe_scaling}
  \caption{\syntheticnote{moe_scaling}GLM-5 MFU (on active parameters) vs.\ GPUs, by expert-parallel size, BF16 and
    FP8. Left of the dotted line: truncated models.}
  \label{fig:moe}
\end{figure}

\paragraph{Expert-group placement.} Keeping every EP group inside a switch (aligned) is the cheapest layout,
but it fits on only 52 of the 66 nodes (an even number per leaf). The question is whether using every node
pays off: more GPUs, but some groups crossing switches. \Cref{fig:expert-placement} (left) puts every run
on one plane, with the same scale on both axes: more GPUs than aligned (x) against more total throughput
than aligned, summed over all GPUs (y). The diagonal is throughput proportional to GPUs; below zero, a run
delivers less in total than aligned. On the same 52 nodes (x = 0), only the node order changes, which gives
the cost of a minority of crossing groups (mixed) and of the worst case (crossed). The run on every node is
worth it only if it lands above zero; with $N$ nodes, that requires a per-GPU cost below $1 - 52/N$,
i.e.\ 19\,\% at 64 nodes. The right panel shows where the extra GPUs go: the gain the extra nodes would
bring if throughput were proportional, the cost of crossing groups (the per-GPU cost of mixed vs.\ aligned
on the same nodes, applied to every node: an estimate) and the rest (larger scale, more simultaneous
crossings: the remainder, an estimate too).

\begin{figure}[h]
  \centering\fig{expert_placement}
  \caption{GLM-5 expert-group placement on Rorqual (EP8). Left: more GPUs vs.\ more total throughput than
    the largest aligned run (52 nodes), equal scales; dotted: distance below proportional scaling. Right:
    total throughput of the run on every node vs.\ aligned (= 100), step by step; solid: measured or
    arithmetic, hatched: estimated. \syntheticnote{expert_placement}}
  \label{fig:expert-placement}
\end{figure}

Every rank waits for the slowest group at each all-to-all, so a few crossing groups may cost almost as
much as all of them. \Cref{fig:expert-groups} shows the all-to-all time of each group in the mixed runs.

\begin{figure}[h]
  \centering\fig[0.6\linewidth]{expert_groups}
  \caption{All-to-all time per EP group in the mixed runs, inside one switch vs.\ across switches.
    \syntheticnote{expert_groups}}
  \label{fig:expert-groups}
\end{figure}

\placeholder{cost of mixed and crossed; total-throughput gain of every node over aligned, and whether it
clears break-even.}

\placeholder{GLM-5 configuration (layers, experts, top-$k$, EP and routing modes), output format from the
Megatron pipeline (see \texttt{FORMAT.md}), and results. Per-rank expert load, if logged, would show
routing imbalance.}
